Las pruebas de carga descritas en \cref{SEC:TEST-LOAD} se ejecutaron
con Locust contra el \textit{endpoint} \texttt{/api/v1/} siguiendo los
tres perfiles de carga que se detallan a continuación.

\bigskip\noindent\textbf{Fase 1: ingesta sostenida.}

El primer escenario simula una cabina de escaneo entregando páginas a
través de la \acs{api} de ingesta a tasa constante. El sistema procesó
14\,260 peticiones a 23,8 \textit{RPS} sostenidos sin errores durante
diez minutos, con latencia mediana inferior a un milisegundo.

\begin{figure}[Fase 1: throughput sostenido]%
              {fig:cdetail-phase1-throughput}%
              {Fase 1 (ingesta sostenida): tasa de peticiones a lo
               largo del tiempo.}
  \image{}{}{load_test/graphs/phase1/02_throughput_over_time}
\end{figure}

\bigskip\noindent\textbf{Fase 2: corrección concurrente.}

El segundo escenario simula a un grupo de correctores anotando y
calificando simultáneamente sobre una base de datos poblada. El sistema
procesó cerca de 5\,900 peticiones sin errores, con latencias medianas
entre 14 y 44 milisegundos por operación, satisfaciendo el objetivo de
p95 inferior a 500\,ms para los \textit{endpoints} \acs{crud}.

\begin{figure}[Fase 2: tasa de errores]%
              {fig:cdetail-phase2-errors}%
              {Fase 2 (corrección concurrente): tasa de errores a lo
               largo del escenario.}
  \image{}{}{load_test/graphs/phase2/03_error_rate}
\end{figure}

\bigskip\noindent\textbf{Fase 3: pico de publicación.}

El tercer escenario reproduce la situación de mayor estrés:
300 alumnos consultando y descargando sus exámenes simultáneamente
tras la publicación de calificaciones. Bajo este perfil, los
\textit{endpoints} ligeros mantienen el rendimiento pero
\texttt{student/instance/download} presenta degradación: alrededor del
21\,\% de errores \acs{http} 500 y latencias medianas de 24 segundos,
atribuible al uso del servidor de desarrollo (\texttt{runserver}) en
lugar de un servidor \acs{wsgi} dedicado.

\begin{figure}[Fase 3: latencia frente a usuarios concurrentes]%
              {fig:cdetail-phase3-latency}%
              {Fase 3 (pico de publicación): evolución de la latencia
               con el número de usuarios concurrentes.}
  \image{}{}{load_test/graphs/phase3/05_latency_vs_users}
\end{figure}

\begin{figure}[Fase 3: errores frente a usuarios concurrentes]%
              {fig:cdetail-phase3-errors}%
              {Fase 3 (pico de publicación): aparición de errores con
               el aumento de usuarios concurrentes.}
  \image{}{}{load_test/graphs/phase3/06_errors_vs_users}
\end{figure}
